\section{Discussion}
\label{sec:discussion}

\subsection{Guidance for system integrators}

The measurements translate directly into selection guidance for practitioners
who compose commercial speech services. Above a signal-to-noise ratio of
about 5~dB, all five recognition systems deliver sub-percent to low-percent
error rates on clearly articulated speech, and the choice among them can be
driven by latency, cost, and interface ergonomics rather than accuracy. The
latency data in Table~\ref{tab:latency} then favour the LLM-ASR text-input
systems and the specialised streaming recognisers, which all answer in well
under a second, over the omni-modal path. Between 0 and 5~dB, the
architectural families separate: the degradation gradients in
Table~\ref{tab:slopes} show that every dB of noise costs a specialised streaming recogniser roughly four to five times what it costs an LLM-ASR system in Mandarin and about twice as much in English, so in
intermittently noisy environments such as public spaces, vehicles, or shared
offices, the foundation-model recognisers dominate the accuracy-latency
trade-off. Below 0~dB, no system transcribes English reliably, and the
appropriate response is architectural: capture improvement, multi-microphone
processing, or explicit rejection logic. For telephone-band audio, the
surprising result is that no specialist premium exists; the generalist
systems are as accurate in-band, and the telephone specialist is markedly
worse under reverberation, so its selection can be justified only by its
8~kHz interface requirements and its speed, not by accuracy.

For translation pipelines, the frontend hypothesis dominates all other
choices, and the latency data combine with it into a compact decision rule.
A cascade assembled from the fastest recognition tier and the text
translation engine completes in about 1.1 seconds per sentence at the mean,
roughly twice as fast as the omni end-to-end path, while the streaming
translation service completes in about half a second and is therefore the
natural choice when latency dominates; when noise robustness dominates, the
same streaming service and the LLM-frontend cascade lead the field, so the
two objectives point to the same shortlist. Swapping the recognition stage of a cascade from the specialised
streaming system to the LLM recogniser is worth 28 BLEU points at 0~dB on an
identical translation backend, which exceeds any reasonable expected gain
from changing the translation engine itself. The linear pass-through rule
quantified in Section~\ref{sec:results:cac}, at roughly half a BLEU point
per percentage point of recognition error, allows an integrator to rank
candidate frontends with recognition-only measurements, which are cheaper to
run and easier to interpret than full-pipeline evaluations, and to forecast
translation quality degradation from a measured or contracted recognition
error rate. The end-to-end versus cascade decision should rest on measured
degradation behaviour rather than architectural intuition: the omni-modal
end-to-end path evaluated here is simultaneously the slowest and, under
severe noise, the most fragile of the four translation designs, while a
dedicated streaming translation service achieves both the smallest relative
loss at 0~dB and sub-second latency.

\subsection{Why foundation-model recognisers degrade more slowly}

The robustness advantage of the LLM-ASR systems echoes the finding of
Radford et al.~\cite{radford2023robust} that recognition models trained on
very large and weakly curated corpora acquire substantial noise tolerance
without explicit multi-condition training. Two of our observations sharpen
this picture for API consumers. First, the advantage is strongly
condition-dependent: it is large under energetic babble noise and negligible
under codec, packet-loss, clipping, and tempo distortions, which barely
affect any system. This pattern is consistent with robustness arising from
exposure to diverse recording conditions during pre-training rather than
from any explicit enhancement front-end. Second, the advantage is
language-asymmetric, with Chinese recognition nearly immune at 0~dB while English recognition of the same systems fails at 37 to 50 percent. Providers
disclose neither training mixture nor per-language capacity, so the asymmetry
can only be documented, not explained, at the black-box level; the
documentation itself is actionable, since it implies that multilingual
deployments should be derated per language rather than per provider. It also implies a concrete procurement test: a
single-language noise probe, one hundred sentences of the target language
at 0~dB babble, is enough to place a candidate service on the robustness
spectrum documented here, and repeating the probe per deployment language
guards against extrapolating a strong Mandarin result to a weak English
result or the reverse. The direction decomposition of
Section~\ref{sec:results:cascade} extends the same caution to translation
deployments, where the two directions of a bilingual service degraded by
strikingly different amounts under identical noise.

\subsection{Cascade versus end-to-end revisited}

The classical comparison of Bentivogli et al.~\cite{bentivogli2018cascade}
found end-to-end and cascade systems competitive on clean speech with
different error structures. Our measurements extend that comparison to
degraded channels in the black-box setting and qualify it in two ways. The
end-to-end advantage in absolute quality against an independent reference is
real but partly a family artefact, since the omni system that translates
directly also shares a generation family with the strongest reference
producer, and the advantage inverts under the secondary reference, as
Section~\ref{sec:results:dual} shows. More importantly, robustness is not an
architectural property but a property of the acoustic frontend wherever it
lives. The end-to-end service with the specialised internal frontend is the
most resilient path we measured, while the end-to-end omni path is the most
fragile, and the same frontend ordering that explains cascade behaviour also
explains end-to-end behaviour. For a consumer, the practical question is not
which architecture to prefer but where the provider placed its best acoustic
model, and the answer is discoverable only by measurement of the kind this
benchmark performs.

\subsection{Methodological implications for black-box evaluation}

Three methodological observations generalise beyond the specific services
tested. The dual-reference protocol of Section~\ref{sec:results:dual}
demonstrates that machine-generated evaluation references can reorder
systems in a way that tracks reference-generator affinity rather than
translation quality, and that degradation-relative conclusions are the
stable currency of black-box comparison. The synthesis-based corpus design
gives black-box studies an exact transcript by construction, which matters
because at sub-percent clean error rates, transcription uncertainty in
recorded corpora would be comparable to the effects being measured. Finally,
the refusal behaviour documented in Section~\ref{sec:results:main}, where
one system returns empty output on a small but systematic subset of severely
degraded utterances, argues that black-box benchmarks should report response
completeness alongside error rates, since a missing transcript and a wrong
transcript have different downstream failure modes.

\subsection{A worked example of the pass-through rule}

The pass-through rule is most useful when stated as an explicit calculation.
Suppose an integrator operates a cascade built on the LLM recogniser in an
environment that occasionally degrades to 10~dB, at which the recogniser
errs at 0.9 percent, and contracts a translation stage equivalent to the one
measured here. The rule predicts a translation loss of roughly 0.47 times
0.9, that is 0.4 BLEU points, which is below the clean-condition agreement
between the two reference systems and therefore negligible. If the same
deployment used the specialised streaming recogniser, whose 10~dB error
rate is 2.8 percent, the predicted loss is 1.9 BLEU points, still modest,
but the margin to the 0~dB regime, where the rule predicts 27 to 38 points
of loss for the two frontends respectively, is thin. The rule therefore
both quantifies the mild regime, where no architectural action is needed,
and identifies the severity threshold beyond which the frontend choice
becomes the dominant design decision. The same arithmetic applied to the
measured error rates of a candidate system at the actual operating
signal-to-noise distribution of a deployment yields an expected translation
quality distribution without running the translation stage at all.

\subsection{Implications for evaluation practice and future work}

Beyond the specific findings, the benchmark design carries lessons for
future evaluations of commercial speech intelligence. A benchmark intended
to track black-box services over time should log every raw response rather
than only aggregate scores, because behavioural changes such as the refusal
pattern documented here appear as distributional shifts before they appear
in means. It should treat reference construction as a first-class
experimental variable, since Section~\ref{sec:results:dual} shows that the
reference generator can reorder systems. It should separate degradation
gradients, which generalise across severity, from point estimates, which do
not, and it should report response completeness alongside accuracy. The
present study points to four directions for follow-up work: replication
across providers to separate family effects from provider effects, extension
to spontaneous and accented speech, sweep-based characterisation of the
currently fixed-severity conditions, and human-reference scoring of the
translation stage to anchor the machine-reference comparisons. As foundation
model services begin to expose streaming speech interfaces with embedded
recognition, the measurement approach demonstrated here, treating each
service as an opaque function from degraded audio to text and scoring the
complete pipeline end to end, appears to be the only viable methodology for
the comparative robustness questions that practitioners actually face.

\section{Limitations}
\label{sec:limitations}

The benchmark uses synthetic speech from four neural voices. Synthetic
speech lacks the disfluencies, accents, emotional variation, and channel
noise of spontaneous human recordings, and although the exact-transcript
property is decisive for measurement at sub-percent error floors, clean
error rates measured on synthetic speech should not be read as population
estimates for human speech; the degradation deltas, which are the object of
this study, are less sensitive to this concern because they are within-system
contrasts. All systems come from a single provider, so architectural family
effects and provider effects are confounded; the family-level regularities
reported here are consistent with published training-scale
explanations~\cite{radford2023robust}, but cross-provider replication is
required before they can be attributed to architecture rather than to one
provider's engineering. The babble noise is itself synthesised from
disjoint neural voices, which makes it speech-shaped and fully reproducible
but smoother than real multi-talker cocktail-party interference, which
includes semantic competition and stronger nonstationarity. Severity levels
for the non-noise degradations were fixed at deployed-but-aggressive values
rather than swept, so only the noise family is characterised by a gradient.
Translation quality is measured with BLEU and chrF against machine-generated
references, with the family-affinity bias that
Section~\ref{sec:results:dual} quantifies; human-reference scoring would
strengthen absolute claims. Finally, commercial APIs are moving targets: the
measurements document the October 2026 versions of the listed services, and
one of the systems could not complete its full condition grid before its
free-tier quota was exhausted, which is itself a practically relevant datum
for anyone planning a similar measurement on complimentary quotas.
